
\vspace*{-1mm}
\begin{ledgroupsized}[r]{120mm}
\footnotesize
\pstart
\noindent\textbf{Überlieferung:}
\pend
\end{ledgroupsized}
\begin{ledgroupsized}[r]{114mm}
\footnotesize
\pstart
\parindent -6mm
\makebox[6mm][l]{\textit{L}}%
Zeichnungen:
LH~XXXVII~4, Bl.~3.
Ein Zettel (ca.~5,5~x~18,5~cm.); oberer und unterer Rand beschnitten;
Papiererhaltungsmaßnahmen.
Eine Seite auf Bl.~\href{\Scan{LH_37_04_003v}}%
{3~v\textsuperscript{o}}.
Bl.~\href{\Scan{LH_37_04_003r}}%
{3~r\textsuperscript{o}}
überliefert einen Text geometrischen Inhalts mit dem Incipit "Sumtis in circulo \textit{F} cujus centrum\;\textit{A}" (das Stück wird in einem späteren Band der Siebenten Reihe erscheinen). 
Ein auf Bl.~3~v\textsuperscript{o} befindlicher, zu letzterem Stück gehöriger Diagrammentwurf wird nicht wiedergegeben.
\pend
\end{ledgroupsized}
\vspace{5mm}
\begin{ledgroup}
\footnotesize
\pstart
\noindent
\textbf{Datierungsgründe:}
Die Zeichnungen in N.\;\ref{60221-B} lassen sich als Abbildungen verschiedener Stoßmöglichkeiten zwischen Kugeln auf einem Billardtisch deuten.
Um die Ziele, d.\,i. auf dem Tisch freistehende Bögen, zu treffen (zu dieser Variante des Billardspiels siehe bspw. das Titelbild von \textsc{C.\;Cotton} \bracks{anonym}, \cite{02159}\textit{The Compleat Gamester}, London 1674) sollen 
verschiedene Eigenschaften des Stoßes eingesetzt werden.
So hebt Leibniz z.\,B. in den 
\bracks{\textit{Fig.~2}}, \bracks{\textit{Fig.~6}} und \bracks{\textit{Fig.~8}} ("Spiel über Bande") 
das Theorem der Gleichheit der Reflexionswinkel,
in den \bracks{\textit{Fig.~3}}, \bracks{\textit{Fig.~4}}, \bracks{\textit{Fig.~5}} und \bracks{\textit{Fig.~7}} indessen
die Phänomene des Mehrkörperstoßes hervor, mit deren Hilfe eine Kugel indirekt ins Ziel gestoßen werden kann.
\pend
\pstart
Das Fehlen jeglichen Kommentars oder Begleittextes zu den Zeichnungen, wie auch eines Wasserzeichens im Träger, stellt ein Hindernis zur Datierung von N.\;\ref{60221-B} dar.
Auch liegt für die noch nicht edierte Aufzeichnung geometrischen Inhalts auf Bl.~3~r\textsuperscript{o}, die Leibniz möglicherweise zeitnah zu N.\;\ref{60221-B} verfasste, derzeit keine verlässliche Datierung (durch die Siebente Reihe) vor.
\pend
\pstart
Zwar besteht zwischen N.\;\ref{60221-B} und einem um 1686 bis Oktober 1687 entstandenen Konzept (\cite{02093}"Principia universalia ad concursus directos ac obliquos determinandos": \textit{LSB} VIII,~3 N.\;66, bes. S.\;715\,f.) eine gewisse inhaltliche Ähnlicheit, da Leibniz auch dort sich mit den "leges Tabulae Globivolvae, quam Germani trucktafel, Galli billard vocant", also mit den theoretischen Grundlagen des Billardspiels, beschäftigt.
Die Ähnlichkeit wird jedoch nicht als hinreichend belastbarer Grund erachtet, die Entstehung von N.\;\ref{60221-B} um dieselbe Zeit anzunehmen und dadurch einer noch nicht erfolgten Datierung der Aufzeichnung auf Bl.~3~r\textsuperscript{o} vorzugreifen.
Plausibel erscheint nach aktuellem Wissensstand und unter den genannten Vorbehalten lediglich die Vermutung, dass es sich um ein Stück aus Leibnizens Hannoveraner Zeit handele.
\pend 
\end{ledgroup}
\selectlanguage{latin}		%
\vspace{5.5mm}
\pstart
\noindent
\href{\Scan{LH_37_04_003v}}%
{\lbrack3~v\textsuperscript{o}\rbrack} 
\pend
\nobreak
\begin{edfigure}[-.4em][0pt]
\hfill
\raisebox{2ex}{
\includegraphics[scale=1]{%
gesamttex/edit_VIII-4/images/LH_37_04_003v_d1.pdf}}
\hfill%
\includegraphics[scale=.25]{%
gesamttex/edit_VIII-4/images/LH_37_04_003v_d1_faks_SW.jpg}		%
\hfill
\edcaption{%
\hspace*{5mm}%
\lbrack\textit{Fig.~1}\rbrack 
\hspace*{40mm}%
\lbrack\textit{Fig.~1a (Bl.~3~v\textsuperscript{o})}\rbrack}
\end{edfigure}
\clearpage
\begin{edfigure}[0ex][0pt]
\hfill
\raisebox{2.5ex}{%
\includegraphics[scale=1]{%
gesamttex/edit_VIII-4/images/LH_37_04_003v_d2.pdf}}
\hfill%
\includegraphics[scale=.25]{%
gesamttex/edit_VIII-4/images/LH_37_04_003v_d2_faks_SW.jpg}	%
\hfill
\edcaption{\lbrack\textit{Fig.~2}\rbrack 
\hspace*{35mm}%
\lbrack\textit{Fig.~2a (Bl.~3~v\textsuperscript{o})}\rbrack}
\end{edfigure}
\begin{edfigure}[4ex][0pt]
\hspace*{20mm}%
\hfill
\raisebox{6ex}{%
\includegraphics[scale=1]{%
gesamttex/edit_VIII-4/images/LH_37_04_003v_d3.pdf}}
\hfill\hfill
\includegraphics[scale=.25]{%
gesamttex/edit_VIII-4/images/LH_37_04_003v_d3+4+5_faks_SW.jpg}	%
\hfill
\edcaption{%
\edtext{%
\hspace*{5mm}%
\lbrack\textit{Fig.~3}\rbrack 
\hspace*{35mm}%
\lbrack\textit{Fig.~3a (Bl.~3~v\textsuperscript{o})}\rbrack%
}{%
\lemma{\bracks{\textit{Fig.~3}} -- \bracks{\textit{Fig.~5}}}%
\Cfootnote[nonum]{%
Leibniz hat im reproduzierten Bereich der Handschrift \bracks{\textit{Fig.~3a}} vermutlich drei verschiedene Stoßkonfigurationen veranschaulicht, bei denen die Kugeln jeweils ins Ziel (das editorisch mit \bracks{\textit{A}} bezeichnete Element) treffen.
Der Übersichtlichkeit halber werden die drei Fälle unter Wiederholung des Elements \bracks{\textit{A}} gesondert wiedergegeben.}}}
\end{edfigure}
\begin{edfigure}[2.5ex][0pt]
\hfill%
\raisebox{1.5ex}{%
\includegraphics[scale=1]{%
gesamttex/edit_VIII-4/images/LH_37_04_003v_d4.pdf}}
\hfill%
\includegraphics[scale=1]{%
gesamttex/edit_VIII-4/images/LH_37_04_003v_d5.pdf}
\hfill
\edcaption[0pt]{\lbrack\textit{Fig.~4}\rbrack
\hspace*{45mm}%
\lbrack\textit{Fig.~5}\rbrack%
\hspace*{5mm}}
\end{edfigure}
\clearpage
\begin{edfigure}[5ex][0pt]
\hfill
\raisebox{2.ex}{%
\includegraphics[scale=1]{%
gesamttex/edit_VIII-4/images/LH_37_04_003v_d6.pdf}}
\hfill%
\includegraphics[scale=.25]{%
gesamttex/edit_VIII-4/images/LH_37_04_003v_d6_faks_SW.jpg}	%
\hfill
\edcaption{\lbrack\textit{Fig.~6}\rbrack 
\hspace*{40mm}%
\lbrack\textit{Fig.~6a (Bl.~3~v\textsuperscript{o})}\rbrack%
\hspace*{5mm}}
\end{edfigure}
\begin{edfigure}[5ex][0pt]
\hfill
\raisebox{5ex}{%
\includegraphics[scale=.8]{%
gesamttex/edit_VIII-4/images/LH_37_04_003v_d7.pdf}}
\hfill%
\includegraphics[scale=.25]{%
gesamttex/edit_VIII-4/images/LH_37_04_003v_d7_faks_SW.jpg}	%
\hfill
\edcaption{\lbrack\textit{Fig.~7, gestrichen}\rbrack 
\hspace*{30mm}%
\lbrack\textit{Fig.~7a (Bl.~3~v\textsuperscript{o})}\rbrack%
\hspace*{10mm}}
\end{edfigure}
\begin{edfigure}[5ex][0pt]
\hfill\hspace*{5mm}%
\raisebox{6ex}{%
\includegraphics[scale=.8]{%
gesamttex/edit_VIII-4/images/LH_37_04_003v_d8.pdf}}
\hfill\hspace*{10mm}%
\includegraphics[scale=.25]{%
gesamttex/edit_VIII-4/images/LH_37_04_003v_d8_faks_SW.jpg}	%
\hfill
\edcaption{\lbrack\textit{Fig.~8, gestrichen}\rbrack 
\hspace*{35mm}%
\lbrack\textit{Fig.~8a (Bl.~3~v\textsuperscript{o})}\rbrack%
\hspace*{12mm}}
\end{edfigure}
